\documentclass[12pt]{article}
\usepackage[T1]{fontenc}
\usepackage[utf8]{inputenc}
\usepackage{lmodern}
\usepackage{textcomp}
\usepackage{enumitem}
\usepackage{geometry}
\usepackage{graphicx}
\usepackage{amsmath, amssymb}
\geometry{margin=1in}
\begin{document}
\begin{center}
Astronomy - Graduate Level Assessment 8 \
Total Marks: 40 \
Answer all questions.
\end{center}

\section*{Multiple Choice (10 marks)}
\begin{enumerate}[label=\arabic*.]
\item The lunar maria are:
\begin{enumerate}[label=(\alph*)]
    \item ancient heavily cratered highlands
    \item dark lavas inside volcanic calderas
    \item dark lavas filling older impact basins
    \item the bright regions on the Moon
\end{enumerate}
\item What effect or effects would be most significant if the Moon's orbital plane were exactly the same as the ecliptic plane?
\begin{enumerate}[label=(\alph*)]
    \item Solar eclipses would be much rarer.
    \item Solar eclipses would last much longer.
    \item Solar eclipses would be much more frequent.
    \item Solar eclipses would not last as long.
\end{enumerate}
\item How do we know how old the Earth is?
\begin{enumerate}[label=(\alph*)]
    \item From the layering of materials within the Earth.
    \item From fossils of ancient life.
    \item From the cratering history of Earth's surface.
    \item From radioactive dating of rocks and meteorites.
\end{enumerate}
\item The constellation ... is a bright W-shaped constellation in the northern sky.
\begin{enumerate}[label=(\alph*)]
    \item Centaurus
    \item Cygnus
    \item Cassiopeia
    \item Cepheus
\end{enumerate}
\item Why is the Mars Exploration Rover Spirit currently tilted towards the north?
\begin{enumerate}[label=(\alph*)]
    \item Because it's climbing up a big hill.
    \item Because it's in the southern hemisphere where it is winter now.
    \item Because it's in the northern hemisphere where it is winter now.
    \item Because one of its wheels broke.
\end{enumerate}
\end{enumerate}

\section*{True / False (10 marks)}
\textit{Answer True or False}
\begin{enumerate}[label=\arabic*.]
\setcounter{enumi}{5}
\item True or False: The Moon's rotational and orbital periods are equal, resulting in the same face always facing the Earth.
\item True or False: Meteorites with high metal content are often fragments of large differentiated asteroids that have been shattered by collisions.
\item True or False: A neutral atom always has equal numbers of neutrons and protons, ensuring stability and balance in the atomic structure.
\item True or False: The telescope was invented by Galileo in 1709, marking the beginning of modern astronomical study.
\item True or False: Meteorites always exhibit a fusion crust due to their passage through the Earth's atmosphere.
\end{enumerate}

\section*{Long Form Response (20 marks)}
\textit{Answer each question with a clear and complete explanation.}
\begin{enumerate}[label=\arabic*.]
\setcounter{enumi}{10}
\item Discuss the significance of Barnard's Star in the context of stellar proximity to Earth, and analyze its characteristics compared to Proxima Centauri and Alpha Centauri.
\item Discuss the rationale behind Ptolemy's use of 'circles upon circles' in his geocentric model of the universe, particularly in relation to the observed retrograde motion of planets.
\end{enumerate}

\vfill
\noindent\textit{End of Paper}
\end{document}